\documentclass[11pt,a4paper]{article}
\usepackage[margin=1in]{geometry}
\usepackage[T1]{fontenc}
\usepackage{lmodern}
\usepackage{booktabs}
\usepackage{longtable}
\usepackage{array}
\usepackage{xcolor}
\usepackage{hyperref}
\usepackage{enumitem}
\hypersetup{colorlinks=true,linkcolor=blue,urlcolor=blue,pdftitle={Technical Progress Report}}
\newcommand{\project}{Nepali ASR/TTS: Syllable Tokenizer and Balanced Corpus Builder}
\newcommand{\status}[1]{\textcolor{blue}{\textbf{#1}}}
\begin{document}
\begin{center}
{\LARGE\bfseries Technical Progress Report\par}
\vspace{0.35em}
{\large \project\par}
\vspace{0.5em}
\textbf{Report date:} August 5, 2026 \qquad
\textbf{Overall status:} \status{Operational and validated on representative fixtures}
\end{center}
\section{Purpose}
This report records the technical progress of the Nepali ASR corpus-building workflow.
It is designed for end-of-day updates so that decisions, validation evidence, risks, and next steps remain traceable.
\section{Current System Snapshot}
\begin{longtable}{>{\raggedright\arraybackslash}p{0.28\textwidth}>{\raggedright\arraybackslash}p{0.64\textwidth}}
\toprule
\textbf{Area} & \textbf{Current state} \\
\midrule
Target corpus & 50,000 selected Nepali ASR sentences, produced in ten batches of 5,000. \\
Primary inputs & Nepali-Text-Corpus Parquet shards, optional \texttt{compiled.txt}, and optional \texttt{Source\_book.txt}. \\
Metadata & Tense, polarity, gender, and sector labels are assigned through configurable rule-based annotation. \\
Selection & Metadata-stratified, syllable-aware greedy selection; prior selected IDs are excluded through \texttt{corpus\_state.json}. \\
Final artifact & \texttt{asr\_corpus/corpus\_50k.jsonl} under the tokenizer dataset directory. \\
Statistics & \texttt{asr\_corpus/reports/} and cumulative \texttt{asr\_corpus/corpus\_state.json}. \\
\bottomrule
\end{longtable}
\section{Workflow Architecture}
\begin{enumerate}[leftmargin=1.7em]
\item \textbf{Extract:} Stream source files, split text into sentences, filter quality, deduplicate normalized text, and write chunked JSONL pool files.
\item \textbf{Annotate:} Add tense, polarity, gender, and sector labels to every pool chunk.
\item \textbf{Select:} Stream annotated chunks, retain a bounded candidate reservoir for each metadata cell, and select the next balanced batch.
\item \textbf{Analyze:} Write a batch report with metadata distributions and syllable statistics.
\item \textbf{Merge:} Deduplicate selected batch records and write the final corpus JSONL file.
\end{enumerate}
\section{Completed Work}
\subsection{Corpus and Distribution Status}
The existing output was checked for consistency. The ten batch files contain 50,000 records with 50,000 unique pool IDs.
The merged corpus also contains 50,000 records. The cumulative state records 10 completed batches, 841 unique syllables, and a cumulative syllable coefficient of variation (CV) of 2.972.
\begin{longtable}{lrr}
\toprule
\textbf{Metadata axis} & \textbf{Label} & \textbf{Count} \\
\midrule
\endhead
Tense & Future & 8,437 \\
& Mixed & 15,595 \\
& Past & 13,007 \\
& Present & 12,961 \\
\addlinespace
Polarity & Negative & 16,667 \\
& Neutral & 18,863 \\
& Positive & 14,470 \\
\addlinespace
Gender & Feminine & 15,324 \\
& Masculine & 15,183 \\
& Neutral & 19,493 \\
\bottomrule
\end{longtable}
Sector-level counts and per-batch distribution details are in the report directory.
Download it with \texttt{corpus\_50k.jsonl} and \texttt{corpus\_state.json} after a cloud run.
\subsection{Memory-Safety Improvement}
\textbf{Problem.} The original selection path read every \texttt{pool\_chunk\_*.jsonl} file into one Python list.
On a large cloud run this attempted to materialize all 2,031 pool chunks and created an unnecessary RAM demand.
\textbf{Implemented solution.} Selection now streams every pool chunk and uses reservoir sampling to retain only a bounded number of candidates per metadata cell.
The default is 2,000 candidates per cell, configured with \texttt{--max-candidates-per-cell}; the complete pool is no longer held in memory during selection.
In multi-worker mode this is a global per-cell cap: the allowance is divided between workers and chunk files are interleaved, preventing the cap from multiplying by the worker count.
\subsection{Cloud CPU Scaling Improvement}
The pipeline accepts \texttt{--workers} and defaults to all CPU cores visible to Python.
The recommended cloud invocation is:
\begin{verbatim}
CLOUD_WORKERS="$(nproc)"
python -m dataset_builder.pipeline run-batch \
  --batch-id 1 --target-size 5000 --max-corpus 0 \
  --max-chunks 0 --workers "$CLOUD_WORKERS"
\end{verbatim}
Parallel Parquet extraction, hash-partitioned deduplication, chunk writing, annotation, streaming, and selection use this worker allocation.
The previous 32-worker cap in the extraction merge phase was removed. A one-worker path is available for low-memory or interactive-notebook execution.
\subsection{Workflow Reliability Improvement}
Partial annotation can bias metadata selection. The pipeline now checks pool chunks without materializing the pool and stops with a clear error if annotation is incomplete.
Annotation processes sequentially when one worker is requested, avoiding a macOS/interpreter process-spawn issue in interactive sessions.
\section{Validation Evidence}
\begin{longtable}{>{\raggedright\arraybackslash}p{0.30\textwidth}>{\raggedright\arraybackslash}p{0.58\textwidth}}
\toprule
\textbf{Check} & \textbf{Result} \\
\midrule
Python syntax compilation & Passed for all scripts under \texttt{scripts/}. \\
2,031-chunk memory fixture & Passed with a cap of 25 candidates per metadata cell; selected 10 unique records without loading the entire pool. \\
Multi-worker selection fixture & Passed with 4 workers over 2,031 chunks; the merged reservoir held 25 candidates and selected 10 unique records. \\
Cloud Batch 1 production run & Completed with 32 workers: 51,431,998 candidates were extracted into 1,056 chunks, annotated, and used to select a balanced 5,000-record batch. \\
Cloud Batch 2 incident & The streaming process pool was terminated during its initial implementation because the 2,000-per-cell limit was applied independently in every worker. The global-cap fix is validated locally; cloud retry is pending. \\
End-to-end fixture & Extract, annotate, streaming selection, report generation, and state writing completed successfully. \\
Parallel Parquet fixture & Passed with 3 accepted records and 3 unique output IDs. \\
Existing corpus consistency & Passed: 10 batches, 50,000 selected records, 50,000 unique IDs, and 50,000 merged records. \\
\bottomrule
\end{longtable}
\section{Known Constraints and Operating Guidance}
\begin{itemize}[leftmargin=1.7em]
\item The full source corpus is large. Download only the final corpus, cumulative state, and reports unless a resume is required.
\item Full-core operation increases throughput, while each worker keeps a bounded reservoir. Reduce \texttt{--max-candidates-per-cell} or use \texttt{--workers 1} if RAM is constrained.
\item Distribution reports describe the selected corpus; they do not replace qualitative ASR-data review or audio/text alignment validation.
\item Record the first full production-scale cloud run with its duration, worker count, input size, output checksum, and any observed failures.
\item If a batch fails after a prior batch completes, do not reset the corpus. Pull the current code and rerun only the failed batch; the state file preserves previous selections.
\end{itemize}
\section{Daily Work Log}
\begin{longtable}{>{\raggedright\arraybackslash}p{0.15\textwidth}>{\raggedright\arraybackslash}p{0.22\textwidth}>{\raggedright\arraybackslash}p{0.51\textwidth}}
\toprule
\textbf{Date} & \textbf{Status} & \textbf{Progress and evidence} \\
\midrule
\endhead
2026-08-05 & In progress & Updated the README and pipeline instructions; replaced full-pool selection loading with bounded streaming reservoirs; enabled full cloud-core scaling; validated single-worker and multi-worker fixtures; confirmed the existing 50k corpus consistency; pushed commit \texttt{41637d2}. A 32-worker cloud run completed Batch 1 from 51,431,998 candidates (1,056 chunks) and selected 5,000 balanced records. Batch 2 exposed a per-worker reservoir memory multiplication; the global-cap fix was implemented and validated locally. Next: pull the fix on cloud and rerun Batch 2 only. \\
\addlinespace
\rule{0pt}{3.6em}\textit{YYYY-MM-DD} & \textit{Planned, active, or complete} & \textit{State the goal, files or commands changed, measured result, validation performed, blockers, and the next action.} \\
\addlinespace
\rule{0pt}{3.6em}\textit{YYYY-MM-DD} & \textit{Planned, active, or complete} & \textit{State the goal, files or commands changed, measured result, validation performed, blockers, and the next action.} \\
\bottomrule
\end{longtable}
\section{End-of-Day Update Procedure}
At the end of each workday, provide the day's commands, outputs, errors, decisions, and completed tasks. Then request:
\begin{quote}
\textit{Update \texttt{docs/technical\_progress\_report.tex} for today's work. Include completed work, measured results, validation, blockers, next steps, and any changes to the corpus or distribution statistics. Compile and verify the PDF afterwards.}
\end{quote}
The update should add a Daily Work Log row and revise the relevant technical sections whenever the design, validation evidence, risks, or results change.
\end{document}
